\documentclass[10pt,a4paper]{amsart}
\usepackage[margin=19mm]{geometry}
\usepackage{amsmath,amssymb,mathtools}
\usepackage[scr=rsfs,cal=euler]{mathalfa}
\usepackage{xfrac}
\usepackage[unicode,hidelinks,bookmarks=false]{hyperref}
\newcommand{\ie}{i.e.\ }
\newcommand{\cf}{cf.\ }
\newcommand{\resp}{respectively\ }
\newcommand{\Subsection}{\subsection}
\providecommand{\nobreakdash}{\nobreak\hbox{-}}
\setlength{\emergencystretch}{2em}
\multlinegap0pt
\usepackage{siunitx}
\usepackage{siunitx}
\usepackage{siunitx}
\usepackage{siunitx}
\usepackage{amssymb}
\usepackage{siunitx}
\usepackage{xcolor}
\usepackage{tikz}
\usepackage{float}
\usepackage{mathtools}
\usepackage[shortlabels]{enumitem}
\usepackage{tcolorbox}
\newtheorem{lemma}{Lemma}
\newtheorem{theorem}{Theorem}
\newtheorem{claim}{Claim}
\usepackage{cleveref}
\crefname{lemma}{Lemma}{Lemmas}
\crefname{theorem}{Theorem}{Theorems}
\crefname{claim}{Claim}{Claims}
\DeclareMathOperator{\ar}{ar}
\newcommand{\cB}{\mathcal{B}}
\title{Anti-Ramsey numbers for cancellative configurations in $p$-graphs}
\author{Cheng Chi, Long-tu Yuan}
\date{Source: arXiv:2604.21834v2 (CC BY 4.0)}
\begin{document}
\maketitle

\noindent\emph{Source and licence.} Anti-Ramsey numbers for cancellative configurations in $p$-graphs. Version arXiv:2604.21834v2 (CC BY 4.0), distributed by arXiv under the Creative Commons Attribution 4.0 International licence, http://creativecommons.org/licenses/by/4.0/, as recorded in the arXiv licence metadata for this exact version and shown on the abstract page https://arxiv.org/abs/2604.21834v2. Full licence notice: \texttt{licenses/arxiv-2604.21834-CC-BY-4.0.txt}.\par\smallskip

\noindent\textit{Editorial scope.} Counted unit: the article's theorem thm:asymptotic-all-n (source0.tex lines 1588--1593). The article's complete original proof of it is delivered with it (source0.tex lines 1595--2193, one \texttt{proof} environment of 599 lines). No mathematics is written, repaired or re-derived: the statement and the proof are the article's own lines, in the article's own notation, and no retained line is substituted or re-flowed.  Retained uncounted as the source-local context the proof needs: lines 1483--1488.  Nothing was omitted from any retained span, so the package carries no omission record.  No retained line contains a citation, so the wrapper carries no bibliography and no bibliography entry.  No retained line contains a figure environment, an image-inclusion command or drawing code, so no image is delivered and the package declares no asset dependency.  Editorial formatting only, in addition: an AMS \texttt{amsart} preamble that restates the article's own declarations and macros used by the retained lines, namely the environments \texttt{lemma}, \texttt{theorem}, \texttt{claim} on separate counters, together with the standard packages those lines need.  The retained environments print in this document as Lemma 1, Theorem 1, Claim 1 in order of appearance, whereas the article prints the same items with the numbering of its own full document.  The complete native source of the article as distributed in the arXiv e-print for this exact version is bundled unchanged as \texttt{sources/199050/source0.tex}.
\begin{lemma}\label{lem:gallai-rho}
    Let $Q$ be an edge-colored complete graph with no rainbow triangle. Then
    \begin{equation}\label{eq:gallai-rho}
        c(Q)+\rho(Q)\le |V(Q)|-1.
    \end{equation}
\end{lemma}

\begin{theorem}\label{thm:asymptotic-all-n}
    For every $n\ge4$,
    \[
        \ar(n,F_4)\le \frac{5n^2-8n}{21}.
    \]
\end{theorem}

\begin{proof}
    Let $G$ be a rainbow $F_4$-free edge-coloring of $K_n^{(3)}$.

    The proof has two main steps.  We first derive a global upper bound by
    applying \cref{lem:gallai-rho}.
    We then derive a complementary lower bound by isolating the singleton-colored triples and using the structure they impose on the remaining colors.
    Comparing these two bounds yields the desired estimate.

    For every color $\lambda$ in $G$, define its support by
    \[
        V_\lambda=\{x\in V(G)\colon \text{there is an edge } e \text{ with } x\in e \text{ and } c_G(e)=\lambda\}.
    \]
    Put
    \[
        I=\sum_\lambda |V_\lambda|,
    \]
    where the summation ranges over all colors in $G$.
    Thus $I$ counts vertex-color incidences: a color with support size $r$ contributes $r$, not merely one.
    This weighting is useful because a color with large support already gives many units toward the final count.

    For every vertex $v\in V(G)$, let $L_v$ be the edge-colored link graph on vertex set $V(G)\setminus\{v\}$, where
    \[
        c_{L_v}(xy)=c_G(vxy)
        \qquad\text{for all distinct }x,y\in V(G)\setminus\{v\}.
    \]
    A rainbow triangle in $L_v$ would correspond exactly to the three triples of a rainbow $F_4$ containing $v$.
    Since $G$ is rainbow $F_4$-free, $L_v$ contains no rainbow triangle.
    By \cref{lem:gallai-rho},
    \[
        c(L_v)+\rho(L_v)\le n-2.
    \]
    Define
    \[
        \rho=\sum_{v\in V(G)}\rho(L_v).
    \]
    Then
    \begin{equation}\label{eq:I-rho-upper}
        I+\rho
        =
        \sum_{v\in V(G)}c(L_v)+\sum_{v\in V(G)}\rho(L_v)
        \le n(n-2).
    \end{equation}

    Let $F$ be the $3$-graph consisting of all singleton-colored triples in $G$, that is,
    \[
        E(F)=\{e\in E(G)\colon c_G(e)\text{ is used exactly once in }G\}.
    \]
    We first note that $F$ is a PSTS.
    Indeed, suppose that two triples of $F$ share a pair, say $abc,abd\in E(F)$.
    Let their colors be $\delta_1$ and $\delta_2$.  Since both colors are used
    exactly once and the two triples are distinct, $\delta_1\ne\delta_2$.
    The triple $acd$ cannot have color $\delta_1$ or $\delta_2$, since either
    choice would repeat a singleton color.
    If $c_G(acd)\notin\{\delta_1,\delta_2\}$, then $abc,abd,acd$ form a rainbow copy of $F_4$, a contradiction.
    Thus no two triples of
    $F$ share a pair, and $F$ is a PSTS.

    We shall use the forcing observation repeatedly.
    It is the basic local mechanism of the proof: a singleton triple sharing a pair with a non-singleton triple forces the two remaining triples on the same $4$-set to take the non-singleton color.

    \begin{claim}\label{clm:forcing-new}
        Let $xyz$ be a triple of a non-unique color $\alpha$.  Suppose that
        $xyd\in E(F)$.  Then
        \[
            c_G(xzd)=c_G(yzd)=\alpha.
        \]
    \end{claim}

    \begin{proof}
        Let $\delta=c_G(xyd)$.  The color $\delta$ is used exactly once.
        Consider $xzd$.  It cannot have color $\delta$, because this would
        repeat the singleton color of $xyd$.  If $c_G(xzd)\notin
            \{\alpha,\delta\}$, then the three triples $xyz,xyd,xzd$ form a rainbow copy of $F_4$ on the $4$-set $\{x,y,z,d\}$, impossible.
        Hence $c_G(xzd)=\alpha$.  The proof for $yzd$ is identical.
    \end{proof}

    Write $s=e(F)$.
    Let $L(F)$ be the {\it leave graph} of $F$, namely the graph on $V(G)$ whose
    edges are the pairs not contained in any triple of $F$, and write
    \[
        \ell=e(L(F)).
    \]
    The leave graph records precisely the pairs where the forcing claim cannot be initiated by a singleton triple.
    Since $F$ is a partial Steiner triple system, every edge of $F$ covers three distinct vertex-pairs and no pair is covered twice, so
    \begin{equation}\label{eq:leave-upper-new}
        3s+\ell=\binom{n}{2}.
    \end{equation}

    We now turn to the lower bound. We use the structure of $F$ to analyze the non-unique colors according to how they interact with the pairs covered by \(F\).  This will eventually produce a lower bound for \(I+\rho\).

    We partition the non-unique colors into good colors and bad colors.
    A non-unique color $\alpha$ is called {\it good} if every $\alpha$-colored triple
    has all its three pairs covered by $F$.  Otherwise, $\alpha$ is called {\it bad}.

    We first show that every good color has support size at least seven.
    Let $\alpha$ be a good color and take an $\alpha$-colored triple $abc$.
    Since $\alpha$ is good, there exist vertices $d,e,f$ such that $abd,ace,bcf$ are edges of $F$.
    By \cref{clm:forcing-new}, applied to the triples
    $abc$ and $abd,ace,bcf$, respectively, the six triples
    \[
        acd,\ bcd,\ abe,\ bce,\ abf,\ acf
    \]
    all have color $\alpha$.
    Since $\alpha$ is good and $acd$ is $\alpha$-colored, the pair $cd$ is
    covered by $F$; say
    \[
        cdg\in E(F).
    \]
    Applying \cref{clm:forcing-new} to the $\alpha$-colored triple $acd$ and
    the singleton triple $cdg$ shows that $g\in V_\alpha$.
    The vertices $a,b,c,d,e,f,g$ are distinct.  Indeed, the vertices
    $d,e,f$ lie outside $\{a,b,c\}$, since $abc$ has the non-unique color
    $\alpha$ and hence is not in $F$.  The vertices $d,e,f$ are pairwise
    distinct; otherwise two of the singleton triples $abd,ace,bcf$ would
    share a pair.  Finally, $g\notin\{a,b\}$ because $acd$ and $bcd$ have
    color $\alpha$ and hence are not in $F$, while $g\notin\{e,f\}$ because
    $cdg$ would then share a pair with $ace$ or $bcf$; also
    $g\notin\{c,d\}$ because $cdg$ is a triple.  Therefore
    \begin{equation}\label{eq:good-support-new}
        |V_\alpha|\ge 7
    \end{equation}
    for every good color $\alpha$.

    We now deal with the bad colors.  Let $\cB$ be the set of bad colors.
    Recall that a non-unique color $\alpha$ is bad precisely when there exists
    an $\alpha$-colored triple containing at least one edge of the leave graph
    $L(F)$.  For every $\alpha\in\cB$, fix once and for all one such
    $\alpha$-colored triple, and denote it by $T_\alpha$.
    Thus
    \[
        c_G(T_\alpha)=\alpha
        \qquad\text{and}\qquad
        \binom{T_\alpha}{2}\cap E(L(F))\ne\emptyset.
    \]
    We call $T_\alpha$ the fixed {\it witness} triple of the bad color $\alpha$.
    Define
    \[
        h_\alpha=
        \left|\binom{T_\alpha}{2}\cap E(L(F))\right|.
    \]
    Since $T_\alpha$ is a triple containing at least one leave-edge, we have $1\le h_\alpha\le3$.

    For every leave-edge $e\in E(L(F))$, define
    \[
        K_e=\{\alpha\in\cB: e\subseteq T_\alpha\},
        \qquad
        k_e=|K_e|.
    \]
    Thus $K_e$ consists of the bad colors whose fixed witness triple contains
    the leave-edge $e$.  The families $K_e$ are not meant to form a partition
    of $\cB$: a bad color may belong to more than one of them.


    In the next claims, a unit counted by \(I\) means an ordered pair
    \((\lambda,x)\) with \(x\in V_\lambda\), and a unit counted by \(\rho\)
    means an ordered pair \((z,\lambda)\) such that color \(\lambda\)
    contains a monochromatic triangle in the link graph \(L_z\).

    We first estimate the {\it base contribution} of a bad color, namely the
    contribution coming from its fixed witness triple and from the pairs of
    that triple that are covered by \(F\).


    \begin{claim}\label{clm:bad-base-new}
        Let $\alpha\in\cB$, and write its fixed witness triple as $T_\alpha=xyz$.
        The base contribution of
        $\alpha$ is at least
        \begin{equation}\label{eq:bad-base-new}
            3+2(3-h_\alpha)=9-2h_\alpha
        \end{equation}
        units.
    \end{claim}

    \begin{proof}
        The three vertices $x,y,z$ of $T_\alpha$ are all in $V_\alpha$, so
        they give the three $I$-units
        \[
            (\alpha,x),\qquad (\alpha,y),\qquad (\alpha,z).
        \]


        Among the three pairs of $T_\alpha$, exactly $h_\alpha$ are
        leave-edges.  Hence exactly $3-h_\alpha$ pairs of $T_\alpha$ are
        covered by $F$.

        Let $p$ be one covered pair of $T_\alpha=xyz$, say $p=\{x,y\}$, so that $z$ is the unique vertex of $T_\alpha$ not contained in $p$.
        Since $p$ is covered by $F$, there is a vertex $d_p$ such that $xyd_p\in E(F)$.
        The vertex $d_p$ is not in $T_\alpha$, because otherwise
        $xyd_p=T_\alpha$, which would mean that the non-unique color $\alpha$
        is used on a singleton triple.

        By \cref{clm:forcing-new}, applied to the $\alpha$-colored triple
        $xyz$ and the singleton triple $xyd_p$, we have
        \[
            c_G(xzd_p)=c_G(yzd_p)=\alpha.
        \]
        Therefore $d_p\in V_\alpha$, giving the additional $I$-unit $(\alpha,d_p)$.

        Moreover, in the link graph $L_z$, the three graph-edges $xy, xd_p,yd_p$ all have color $\alpha$, since they correspond to the triples
        $xyz,xzd_p,yzd_p$.
        Hence color $\alpha$ contains a monochromatic
        triangle in $L_z$, giving the $\rho$-unit $(z,\alpha)$.

        These units are distinct as $p$ ranges over the covered pairs of
        $T_\alpha$.  For the $I$-units, if two distinct covered pairs
        $p,q\subseteq T_\alpha$ had $d_p=d_q$, then the two singleton triples $p\cup\{d_p\}$ and $q\cup\{d_q\}$ would share a pair, contradicting that $F$ is a partial Steiner triple
        system.  Also, each $d_p$ lies outside $T_\alpha$, so these units are
        distinct from $(\alpha,x),(\alpha,y),(\alpha,z)$.  For the
        $\rho$-units, distinct pairs of the triple $T_\alpha$ have distinct
        remaining vertices: if $T_\alpha=xyz$, then the pairs $xy,xz,yz$
        correspond respectively to the link vertices $z,y,x$.

        Hence each covered pair of $T_\alpha$ gives two new units, one in $I$
        and one in $\rho$.  Since there are $3-h_\alpha$ covered pairs, the
        base contribution is at least
        \[
            3+2(3-h_\alpha)=9-2h_\alpha.
        \]
    \end{proof}

    We now estimate the additional contribution coming from the leave-edges in the fixed witness triples.  Fix a leave-edge
    \[
        e=uv\in E(L(F)).
    \]
    For every $\alpha\in K_e$, recall that $e\subseteq T_\alpha$.  Since
    $T_\alpha$ is a triple, there is a unique vertex $x_{\alpha,e}$ such that
    \[
        T_\alpha=e\cup\{x_{\alpha,e}\}=uvx_{\alpha,e}.
    \]
    If $\alpha,\beta\in K_e$ are distinct, then
    $x_{\alpha,e}\ne x_{\beta,e}$; otherwise the same triple would have the
    two distinct colors $\alpha$ and $\beta$.

    Consider two distinct colors $\alpha,\beta\in K_e$.  On the $4$-set $\{u,v,x_{\alpha,e},x_{\beta,e}\}$, the triples $uvx_{\alpha,e}$ and $uvx_{\beta,e}$ have colors $\alpha$ and
    $\beta$, respectively.
    Hence each of the two remaining triples $ux_{\alpha,e}x_{\beta,e}$, $vx_{\alpha,e}x_{\beta,e}$ has color in $\{\alpha,\beta\}$; otherwise, together with
    $uvx_{\alpha,e}$ and $uvx_{\beta,e}$, it would form a rainbow copy of
    $F_4$.

    For $t\in\{u,v\}$ and distinct $\alpha,\beta\in K_e$, say that
    $\alpha$ wins over $\beta$ at $t$ with respect to $e$ if
    \[
        c_G(tx_{\alpha,e}x_{\beta,e})=\alpha.
    \]
    For every unordered pair $\{\alpha,\beta\}\subseteq K_e$ and every
    $t\in\{u,v\}$, exactly one of $\alpha$ and $\beta$ wins over the other at
    $t$ with respect to $e$.

    For $\alpha\in K_e$ and $t\in\{u,v\}$, define
    \[
        W_t^e(\alpha)
        =
        \{\beta\in K_e\setminus\{\alpha\}:
        \alpha \text{ wins over } \beta \text{ at } t \text{ with respect to } e\}.
    \]
    For every leave-edge $e=uv$, define
    \[
        \mathcal Q_e^I
        =
        \{(\alpha,x_{\beta,e}):
        \alpha\in K_e,\,
        \beta\in W_u^e(\alpha)\cup W_v^e(\alpha)\}
    \]
    and
    \[
        \mathcal Q_e^\rho
        =
        \{(x_{\alpha,e},\alpha):
        \alpha\in K_e,\,
        W_u^e(\alpha)\cap W_v^e(\alpha)\neq\emptyset\}.
    \]
    We first claim that every element of $\mathcal Q_e^I$ is counted by $I$, and every element of $\mathcal Q_e^\rho$ is counted by $\rho$.
    In fact, if $\beta\in W_u^e(\alpha)$, then $c_G(ux_{\alpha,e}x_{\beta,e})=\alpha$, so $x_{\beta,e}\in V_\alpha$ and $(\alpha,x_{\beta,e})$ is counted by $I$.
    The same argument applies if $\beta\in W_v^e(\alpha)$, using
    the triple $vx_{\alpha,e}x_{\beta,e}$.  Thus every element of
    $\mathcal Q_e^I$ is counted by $I$.
    Write $M_e$ for $|\mathcal Q_e^I|+|\mathcal Q_e^\rho|$.
    Now we count $M_e$ in another way.
    \begin{claim}\label{clm:bad-winner-contribution-new}
        We have
        \begin{equation}\label{eq:Me-def}
            M_e=
            \sum_{\alpha\in K_e}
            \left(
            |W_u^e(\alpha)\cup W_v^e(\alpha)|
            +
            \mathbf 1_{W_u^e(\alpha)\cap W_v^e(\alpha)\neq\emptyset}
            \right),
        \end{equation}
        and
        \begin{equation}\label{eq:Me-lower}
            M_e\ge
            \begin{cases}
                0,                    & k_e=0,   \\
                \binom{k_e}{2}+k_e-1, & k_e\ge1.
            \end{cases}
        \end{equation}
    \end{claim}
    \begin{proof}
        For fixed $\alpha$, the vertices $x_{\beta,e}$ with
        $\beta\in W_u^e(\alpha)\cup W_v^e(\alpha)$ are distinct, since the
        triples $uvx_{\beta,e}$ have distinct colors.  For different values
        of $\alpha$, the first coordinate of the ordered pair
        $(\alpha,x_{\beta,e})$ is different.  Hence
        \[
            |\mathcal Q_e^I|
            =
            \sum_{\alpha\in K_e}|W_u^e(\alpha)\cup W_v^e(\alpha)|.
        \]

        Now suppose that $W_u^e(\alpha)\cap W_v^e(\alpha)\neq\emptyset$.
        Choose $\beta\in W_u^e(\alpha)\cap W_v^e(\alpha)$.  Then
        \[
            c_G(ux_{\alpha,e}x_{\beta,e})
            =
            c_G(vx_{\alpha,e}x_{\beta,e})
            =
            \alpha.
        \]
        Together with $c_G(uvx_{\alpha,e})=\alpha$,
        this shows that in the link graph $L_{x_{\alpha,e}}$ the three edges $uv, ux_{\beta,e}, vx_{\beta,e}$ form a monochromatic triangle of color $\alpha$.
        Thus $(x_{\alpha,e},\alpha)$ is counted by $\rho$.
        This proves that every element of $\mathcal Q_e^\rho$ is counted by $\rho$.
        The elements of $\mathcal Q_e^\rho$ are distinct because their color coordinates
        $\alpha$ are distinct.  Therefore \eqref{eq:Me-def} holds.

        It remains to prove \eqref{eq:Me-lower}.  If $k_e=0$, then
        $M_e=0$.
        Assume $k_e\ge1$. At the endpoint $u$, every unordered pair
        $\{\alpha,\beta\}\subseteq K_e$ has exactly one winner.\footnote{Readers may consider that the ``win'' relation defines a tournament.}  Hence
        \[
            \sum_{\alpha\in K_e}|W_u^e(\alpha)|
            =
            \binom{k_e}{2}.
        \]
        This is the baseline contribution.

        At the endpoint $v$, at most one color $\alpha\in K_e$ can satisfy
        $W_v^e(\alpha)=\emptyset$.  Indeed, if two distinct colors
        $\alpha,\beta\in K_e$ both had empty $W_v^e$-sets, then one of them
        would have to win over the other at $v$, contradicting the emptiness
        of its $W_v^e$-set.  Therefore at least $k_e-1$ colors
        $\alpha\in K_e$ satisfy $W_v^e(\alpha)\neq\emptyset$.
        Fix such an $\alpha$.  If
        \[
            W_u^e(\alpha)\cap W_v^e(\alpha)\neq\emptyset,
        \]
        then the indicator term contributes one additional unit beyond
        $|W_u^e(\alpha)|$.  If
        \[
            W_u^e(\alpha)\cap W_v^e(\alpha)=\emptyset,
        \]
        then
        \[
            |W_u^e(\alpha)\cup W_v^e(\alpha)|
            =
            |W_u^e(\alpha)|+|W_v^e(\alpha)|
            \ge |W_u^e(\alpha)|+1.
        \]
        Thus every $\alpha$ with $W_v^e(\alpha)\neq\emptyset$ contributes at
        least one unit beyond the baseline $|W_u^e(\alpha)|$.

        Since this happens for at least $k_e-1$ colors, we obtain
        \[
            M_e
            \ge
            \binom{k_e}{2}+k_e-1.
        \]
    \end{proof}

    The next claim checks that the base units and the leave-edge units are disjoint.

    \begin{claim}\label{clm:bad-disjoint-new}
        The units counted in the base contribution
        \eqref{eq:bad-base-new} and the units in the families
        $\mathcal Q_e^I,\mathcal Q_e^\rho$, over all leave-edges
        $e\in E(L(F))$, may be counted together in $I+\rho$ without
        repetition.
    \end{claim}

    \begin{proof}
        We check repetitions within $I$ and within $\rho$ separately.  Units
        with different color coordinate are automatically distinct.  Hence it
        is enough to fix one bad color $\alpha$ and prove that no unit
        assigned to this fixed $\alpha$ is counted twice.

        Write $T_\alpha=abc$.
        First consider the units counted by $I$.  The base $I$-units are
        \[
            (\alpha,a),\qquad (\alpha,b),\qquad (\alpha,c),
        \]
        together with the units $(\alpha,d_p)$ coming from covered pairs
        $p\subseteq T_\alpha$, where
        \[
            p\cup\{d_p\}\in E(F).
        \]
        For every covered pair $p$, the vertex $d_p$ lies outside
        $T_\alpha$; otherwise $T_\alpha$ itself would be a singleton triple.
        If $p\ne q$ are two covered pairs of $T_\alpha$, then
        $d_p\ne d_q$, because otherwise the two singleton triples $p\cup\{d_p\}$ and $q\cup\{d_q\}$ would share a pair, contradicting that $F$ is a partial Steiner triple
        system.  Thus the base $I$-units assigned to $\alpha$ are distinct.

        Now consider an $I$-unit assigned to $\alpha$ from some leave-edge
        $e\subseteq T_\alpha$.  Write
        \[
            e=uv,
            \qquad
            T_\alpha=e\cup\{x_{\alpha,e}\}.
        \]
        Such a unit has the form $(\alpha,x_{\beta,e})$, where $\beta\in K_e\setminus\{\alpha\}$ and $T_\beta=e\cup\{x_{\beta,e}\}$.

        First, $x_{\beta,e}\notin T_\alpha$.  Indeed, if
        $x_{\beta,e}\in T_\alpha$, then, since $e\subseteq T_\alpha$, we
        would have $T_\beta=T_\alpha$,
        forcing the same triple to have two distinct colors $\beta$ and
        $\alpha$.

        Second, $x_{\beta,e}$ cannot equal any vertex $d_p$ coming from a
        covered pair $p\subseteq T_\alpha$.  After relabelling, suppose
        \[
            T_\alpha=abc,\qquad e=ab,\qquad T_\beta=aby,
        \]
        and suppose that a covered pair of $T_\alpha$, say $ac$, satisfies $acy\in E(F)$.
        Then the three triples $abc,aby,acy$ lie on the same $4$-set and form a copy of $F_4$.
        Their colors are pairwise distinct: $abc$ has color $\alpha$, $aby$ has color $\beta\ne\alpha$, and $acy$ has a singleton color.
        This gives a rainbow copy of $F_4$, impossible.
        The case where the covered pair is $bc$ is identical by relabelling.

        Third, two $I$-units assigned to $\alpha$ from the same leave-edge
        cannot have the same second coordinate.  If
        $x_{\beta,e}=x_{\gamma,e}$ for distinct $\beta,\gamma\in K_e$, then
        \[
            T_\beta=e\cup\{x_{\beta,e}\}
            =
            e\cup\{x_{\gamma,e}\}
            =
            T_\gamma,
        \]
        so the same triple would have two distinct colors.

        Fourth, two $I$-units assigned to $\alpha$ from two distinct
        leave-edges of $T_\alpha$ cannot have the same second coordinate.
        After relabelling, suppose the two leave-edges are $ab$ and $ac$, and
        suppose both produce the same vertex $y$.  Then there exist bad colors
        $\beta,\gamma$ such that
        \[
            T_\beta=aby,
            \qquad
            T_\gamma=acy.
        \]
        The colors $\alpha,\beta,\gamma$ are pairwise distinct: neither
        $\beta$ nor $\gamma$ equals $\alpha$, and $\beta=\gamma$ would force
        the fixed witness triple $T_\beta$ to be both $aby$ and $acy$.
        Hence, $abc,aby,acy$ form a rainbow copy of $F_4$, impossible.

        Therefore all $I$-units assigned to the fixed color $\alpha$ are
        distinct.

        It remains to check the units counted by $\rho$.  For a pair
        $p\subseteq T_\alpha$, let $z_p$ denote the unique vertex of $T_\alpha\setminus p$.
        Thus, if $T_\alpha=abc$, then
        \[
            z_{ab}=c,\qquad z_{ac}=b,\qquad z_{bc}=a.
        \]

        A covered pair $p\subseteq T_\alpha$ contributes the $\rho$-unit $(z_p,\alpha)$,
        because the monochromatic triangle of color $\alpha$ lies in the link graph $L_{z_p}$.
        A contribution from a leave-edge
        $e=p\subseteq T_\alpha$ through the family $\mathcal Q_e^\rho$ also
        has the form $(z_p,\alpha)$,
        since for $p=uv$ and $T_\alpha=uvz_p$, the corresponding
        $\rho$-unit is $(z_p,\alpha)$.

        However, the same pair $p\subseteq T_\alpha$ cannot be both covered by
        $F$ and a leave-edge of $F$.  Hence the base contribution and the
        leave-edge contribution cannot assign the same $\rho$-unit for the
        same pair $p$.  Moreover, distinct pairs of the triple $T_\alpha$ have
        distinct vertices $z_p$.  Therefore all $\rho$-units assigned to
        $\alpha$ are distinct.

        Since the choice of $\alpha$ was arbitrary, and units with different
        color coordinate are distinct, all the counted units may be counted
        together in $I+\rho$ without repetition.
    \end{proof}

    Let $g$ be the number of good colors.  From the contributions of singleton
    colors, good colors, and bad colors, we obtain
    \begin{equation}\label{eq:I-rho-lower-new}
        I+\rho
        \ge
        3s+7g
        +
        \sum_{\alpha\in\cB}(9-2h_\alpha)
        +
        \sum_{e\in E(L(F))}
        \left(\binom{k_e}{2}+k_e-1\right)_+,
    \end{equation}
    where $x_+=\max\{x,0\}$.

    Let $b=|\cB|$ be the number of bad colors.  Then $c(G)=s+g+b$.
    By \eqref{eq:I-rho-lower-new},
    \begin{align}
        7c(G)
         & =
        7s+7g+7b \notag \\
         & \le
        (I+\rho)+4s
        +
        \sum_{\alpha\in\cB}(2h_\alpha-2)
        -
        \sum_{e\in E(L(F))}
        \left(\binom{k_e}{2}+k_e-1\right)_+.
        \label{eq:seven-c-new}
    \end{align}

    We next rewrite the two remaining sums in \eqref{eq:seven-c-new} as a
    single sum over the leave-edges of \(F\).  For a bad color \(\alpha\), the
    fixed witness triple \(T_\alpha\) contains exactly \(h_\alpha\)
    leave-edges.  We therefore distribute the term \(2h_\alpha-2\) equally
    over these \(h_\alpha\) leave-edge incidences.  Thus each incidence
    \(\alpha\in K_e\) receives weight
    \[
        \frac{2h_\alpha-2}{h_\alpha}
        =
        2-\frac2{h_\alpha}.
    \]
    Equivalently, since a fixed bad color \(\alpha\) belongs to \(K_e\) for
    exactly \(h_\alpha\) leave-edges \(e\subseteq T_\alpha\), a double count
    of these weighted incidences gives
    \begin{equation}\label{eq:weighted-incidence-new}
        \sum_{\alpha\in\cB}(2h_\alpha-2)
        =
        \sum_{e\in E(L(F))}
        \sum_{\alpha\in K_e}
        \left(2-\frac2{h_\alpha}\right).
    \end{equation}
    Since \(h_\alpha\in\{1,2,3\}\), every such weight is at most \(4/3\).

    For a fixed leave-edge \(e\), define
    \[
        R_e=
        \sum_{\alpha\in K_e}
        \left(2-\frac2{h_\alpha}\right)
        -
        \left(\binom{k_e}{2}+k_e-1\right)_+,
    \]
    where $x_+=\max\{x,0\}$.
    By \eqref{eq:weighted-incidence-new},
    \begin{equation}
        \sum_{\alpha\in\cB}(2h_\alpha-2)
        -
        \sum_{e\in E(L(F))}
        \left(\binom{k_e}{2}+k_e-1\right)_+
        =
        \sum_{e\in E(L(F))}R_e.
    \end{equation}


    Since $h_\alpha\in\{1,2,3\}$, this weight is at most $4/3$.
    We claim that
    \begin{equation}\label{eq:Re-bound-new}
        R_e\le \frac{4}{3}.
    \end{equation}
    If $k_e=0$, then $R_e=0$.  If $k_e=1$, then $R_e\le4/3$.  If $k_e=2$,
    then
    \[
        R_e\le \frac{4}{3}+\frac{4}{3}-2=\frac23.
    \]
    If $k_e\ge3$, then
    \[
        R_e
        \le
        \frac{4}{3} k_e-\left(\binom{k_e}{2}+k_e-1\right)
        \le \frac{4}{3}.
    \]
    This proves \eqref{eq:Re-bound-new}.

    Combining \eqref{eq:seven-c-new} and \eqref{eq:Re-bound-new}, we get
    \[
        7c(G)\le (I+\rho)+4s+\frac{4}{3}\ell.
    \]
    By \eqref{eq:I-rho-upper} and \eqref{eq:leave-upper-new},
    \[
        7c(G)
        \le
        n(n-2)+4s+\frac{4}{3}\ell
        =
        n(n-2)+\frac{4}{3}\binom n2.
    \]
    Therefore
    \[
        c(G)
        \le
        \frac17\left(n(n-2)+\frac{4}{3}\binom n2\right)
        =
        \frac{5n^2-8n}{21}.
    \]
    Since $G$ was arbitrary, this proves the theorem.
\end{proof}

\end{document}
